%! Solving Exponential Equations with a Common Base — Unit 7, Lesson 7.3 — Warm-Up
\documentclass[10pt]{article}
\usepackage{algebra2-article}
\usepackage{algebra2-boxes}
\newcommand{\TallMath}[1]{$\displaystyle #1\rule[-1.4em]{0pt}{3.2em}$}

\begin{document}

\pageheader{Unit 7, Lesson 7.3}{Warm-Up}
\namedateperiod

\vspace{0.10in}

\begin{notesbox}{Getting Ready: Skills We'll Use Today}
\small Today you \emph{solve} an exponential equation for the first time. Nothing below is new: two of
these items are Unit 6 exponent work, and the third is Lesson 7.0.

\vspace{5pt}
\textbf{1. The same number in a different outfit.} Fill in each missing exponent.

\vspace{4pt}
\hspace*{1.2em} $4 = 2^{\,\blacksquare}$, \; $\blacksquare=$ \blank{1.1cm} \qquad
$32 = 2^{\,\blacksquare}$, \; $\blacksquare=$ \blank{1.1cm} \qquad
$\dfrac{1}{8} = 2^{\,\blacksquare}$, \; $\blacksquare=$ \blank{1.1cm} \qquad
$\sqrt{2} = 2^{\,\blacksquare}$, \; $\blacksquare=$ \blank{1.1cm}

\vspace{5pt}
\hspace*{1.2em} $81 = 3^{\,\blacksquare}$, \; $\blacksquare=$ \blank{1.1cm} \qquad
$\dfrac{1}{27} = 3^{\,\blacksquare}$, \; $\blacksquare=$ \blank{1.1cm} \qquad
$125 = 5^{\,\blacksquare}$, \; $\blacksquare=$ \blank{1.1cm} \qquad
$\sqrt[3]{7} = 7^{\,\blacksquare}$, \; $\blacksquare=$ \blank{1.1cm}

\vspace{0.32in}

\textbf{2. Two rules that look alike and are not.} Evaluate each one \emph{as a number}, then say what
happened to the exponents.

\vspace{4pt}
\hspace*{1.2em} $2^{3}\cdot 2^{4} =$ \blank{1.8cm} $\;=\; 2^{\,\blacksquare}$ where $\blacksquare=$
\blank{1.1cm} \qquad
$\left(2^{3}\right)^{4} =$ \blank{1.8cm} $\;=\; 2^{\,\blacksquare}$ where $\blacksquare=$
\blank{1.1cm}

\vspace{4pt}
\hspace*{1.2em} Multiplying powers of the same base \blank{2.2cm} the exponents; raising a power to a
power \blank{2.2cm} them.

\vspace{0.38in}

\textbf{3. A question about shape} (Lesson 7.0). Complete the table for $y=2^{x}$.

\vspace{3pt}
\begin{center}
\renewcommand{\arraystretch}{1.35}
\begin{tabular}{|c|c|c|c|c|c|}
\hline
$x$ & $-1$ & $0$ & $1$ & $2$ & $3$ \\ \hline
$y=2^{x}$ & \blank{1.3cm} & \blank{1.3cm} & \blank{1.3cm} & \blank{1.3cm} & \blank{1.3cm} \\ \hline
\end{tabular}
\end{center}

\vspace{2pt}
\hspace*{1.2em} Is there any output this function produces \emph{twice} --- two different inputs giving
the same $y$? \blank{2.4cm}

\vspace{4pt}
\hspace*{1.2em} What was it about the \emph{shape} of an exponential graph in Lesson 7.0 that settles
that question? \writeline

\end{notesbox}

\end{document}
